\section*{Hoofdstuk \ref{ch:g6:wholes} --- Gehele getallen}

\begin{solution}{exo:g6:wholes:1}
$3\,047$; \quad $20\,500$; \quad $104$. \quad
$60\,013$: ``zestigduizend dertien''.
\end{solution}

\begin{solution}{exo:g6:wholes:2}
Cijfer van de honderdtallen: $2$; cijfer van de eenheden: $9$; de $4$ telt de
tienduizendtallen. Splitsing:
\[
47\,209 = 4 \times 10\,000 + 7 \times 1000 + 2 \times 100 + 0 \times 10
+ 9 .
\]
\end{solution}

\begin{solution}{exo:g6:wholes:3}
$507 < 570$ (tientallen: $0 < 7$); \quad
$1\,099 > 989$ (vier cijfers verslaan drie); \quad
$23\,456 < 23\,465$ (tientallen: $5 < 6$); \quad
$9\,999 < 10\,000$ (vier cijfers tegen vijf).
\end{solution}

\begin{solution}{exo:g6:wholes:4}
Op het aantal cijfers komen $978$ en $1\,001$ voor de drie getallen van vier
cijfers die met $2$ beginnen; vergelijk bij die drie de tientallen en de
eenheden: $2\,034 < 2\,043 < 2\,304$. De volgorde is:
\[
978 < 1\,001 < 2\,034 < 2\,043 < 2\,304 .
\]
\end{solution}

\begin{solution}{exo:g6:wholes:5}
$348 + 576 = 924$; \quad $805 - 268 = 537$ (controle:
$537 + 268 = 805$); \quad $307 \times 26 = 307 \times 20 + 307 \times 6
= 6\,140 + 1\,842 = 7\,982$.
\end{solution}

\begin{solution}{exo:g6:wholes:6}
$38 + 45 + 12 + 5 = (38 + 12) + (45 + 5) = 50 + 50 = 100$.

$2 \times 83 \times 50 = (2 \times 50) \times 83 = 100 \times 83 =
8\,300$.

$4 \times 78 \times 25 = (4 \times 25) \times 78 = 100 \times 78 =
7\,800$.
\end{solution}

\begin{solution}{exo:g6:wholes:7}
$95 = 7 \times 13 + 4$ (\omterm{def:g3:division:remainder}{quotiënt} $13$, \omterm{def:g3:division:remainder}{rest} $4$).

$230 = 6 \times 38 + 2$ (\omterm{def:g3:division:remainder}{quotiënt} $38$, \omterm{def:g3:division:remainder}{rest} $2$).

$504 = 8 \times 63 + 0$ (\omterm{def:g3:division:remainder}{quotiënt} $63$, \omterm{def:g3:division:remainder}{rest} $0$: $504$ is \omterm{def:g6:wholes:divisible}{deelbaar} door
$8$).
\end{solution}

\begin{solution}{exo:g6:wholes:8}
De gelijkheid is juist, maar de \omterm{def:g3:division:remainder}{rest} $13$ is \emph{niet} kleiner dan de \omterm{def:g5:division:multiple}{deler}
$8$, dus is het niet de deling met \omterm{def:g3:division:remainder}{rest}. Haal nog één keer $8$ uit de \omterm{def:g3:division:remainder}{rest}:
$173 = 8 \times 21 + 5$, met $0 \leq 5 < 8$: \omterm{def:g3:division:remainder}{quotiënt} $21$, \omterm{def:g3:division:remainder}{rest} $5$.
\end{solution}

\begin{solution}{exo:g6:wholes:9}
\emph{1.} $2\,000 = 12 \times 166 + 8$: de boerderij vult $166$ volle dozen en
er blijven $8$ eieren over.

\emph{2.} $150 = 12 \times 12 + 6$: twaalf dozen bevatten maar $144$ eieren,
en dat is te weinig. De bakkerij moet $13$ dozen kopen.
\end{solution}

\begin{solution}{exo:g6:wholes:10}
$100 = 7 \times 14 + 2$: honderd dagen zijn $14$ volle weken en $2$ dagen
erbij. Veertien weken later is het weer dinsdag; twee dagen na een dinsdag is
het \emph{donderdag}.
\end{solution}

\begin{solution}{exo:g6:wholes:11}
De grootst mogelijke \omterm{def:g3:division:remainder}{rest} bij een deling door $9$ is $8$ (de \omterm{def:g3:division:remainder}{rest} moet kleiner
zijn dan de \omterm{def:g5:division:multiple}{deler}). Het deeltal is dus
\[
9 \times 56 + 8 = 504 + 8 = 512 .
\]
\end{solution}

\begin{solution}{exo:g6:wholes:12}
Noem de cijfers $h$ (honderdtallen), $0$ (tientallen) en $e$ (eenheden). We
weten dat $h = 2 \times e$ en $h + 0 + e = 9$, dus $2e + e = 9$, waaruit
$3e = 9$, $e = 3$ en $h = 6$. Het getal is $603$. Controle: $6$ is twee keer
$3$, het cijfer van de tientallen is $0$, en $6 + 0 + 3 = 9$.
\end{solution}

\begin{solution}{pb:g6:wholes:1}
\textbf{1.} Bijvoorbeeld $(1 + 4) + (2 + 3) + 5 = 5 + 5 + 5 = 15$.

\textbf{2.} Elk paar is $11$ waard: $1 + 10$, $2 + 9$, $3 + 8$, $4 + 7$,
$5 + 6$. Dat zijn $5$ paren (tien getallen, twee per paar), dus is de \omterm{def:g1:addition:def}{som}
$5 \times 11 = 55$.

\textbf{3.} Paren van $21$ ($1 + 20$, $2 + 19, \dots$), en het zijn er $10$:
$10 \times 21 = 210$.

\textbf{4.} In elk paar groeit het eerste getal met $1$ terwijl het tweede met
$1$ krimpt, dus verandert het totaal nooit:
$1 + 100 = 2 + 99 = 3 + 98 = \dots = 101$. De honderd getallen vormen $50$
paren (twee getallen elk), met $50 + 51 = 101$ als binnenste paar.

\textbf{5.} $50$ paren van $101$:
\[
1 + 2 + \dots + 100 = 50 \times 101 = 5\,050 .
\]

\textbf{6.} Bij negen getallen heeft het middelste, $5$, geen partner.
Eromheen: $1 + 9 = 2 + 8 = 3 + 7 = 4 + 6 = 10$, vier paren. Totaal:
$4 \times 10 + 5 = 45$.

\textbf{7.} Elke kolom is $1 + 9 = 2 + 8 = \dots = 9 + 1 = 10$ waard, en er
zijn $9$ kolommen: samen zijn de twee regels $9 \times 10 = 90$ waard. Maar de
twee regels zijn \emph{dezelfde} \omterm{def:g1:addition:def}{som}, twee keer opgeschreven, dus is één
exemplaar de \omterm{def:g3:division:half}{helft} waard: $1 + 2 + \dots + 9 = 90 \div 2 = 45$. Nu blijft er
nooit een getal zonder partner --- zijn partner staat er recht onder.

\textbf{8.} Voor $1$ tot $50$: $50$ kolommen van $51$, dubbele \omterm{def:g1:addition:def}{som}
$50 \times 51 = 2\,550$, dus de \omterm{def:g1:addition:def}{som} is $2\,550 \div 2 = 1\,275$. Voor $1$ tot
$200$: $200$ kolommen van $201$, dubbele \omterm{def:g1:addition:def}{som} $40\,200$, \omterm{def:g1:addition:def}{som} $20\,100$.

\textbf{9.} Elk \omterm{def:g3:numbers:evenodd}{even getal} is het dubbele van een getal uit de \omterm{def:g1:addition:def}{som} van Gauss:
$2 = 2 \times 1$, $4 = 2 \times 2$, \dots, $200 = 2 \times 100$. De hele \omterm{def:g1:addition:def}{som}
is dus het dubbele van die van Gauss:
\[
2 + 4 + \dots + 200 = 2 \times 5\,050 = 10\,100 .
\]

\textbf{10.} De getallen van $101$ tot $200$ zijn de getallen van $1$ tot
$200$ zonder de getallen van $1$ tot $100$:
\[
101 + 102 + \dots + 200 = 20\,100 - 5\,050 = 15\,050 .
\]

\textbf{11.} Zet de vrienden op een rij. De eerste geeft $7$ handen (één aan
elk van de anderen). De tweede heeft de eerste al begroet, dus geeft hij $6$
\emph{nieuwe} handen; de derde $5$; en zo verder, tot de zevende er $1$ geeft
en de achtste geen nieuwe meer. Elke handdruk wordt precies één keer geteld:
\[
7 + 6 + 5 + 4 + 3 + 2 + 1 = 28 \text{ handdrukken}
\]
(koppel om het midden: $(7+1) + (6+2) + (5+3) + 4 = 8 + 8 + 8 + 4 = 28$).

\textbf{12.} $1 + 2 + \dots + 12$: twaalf kolommen van $13$, dubbele \omterm{def:g1:addition:def}{som}
$12 \times 13 = 156$, dus $156 \div 2 = 78$ blikjes.

\textbf{13.} De trap van $1 + 2 + 3 + 4$ vakjes, met dezelfde trap op zijn kop
erbij, vult precies een rechthoek van $4$ rijen en $5$ kolommen: elke rij van
de rechthoek is één kolom van de truc met twee regels (een trede van de
stijgende trap, aangevuld door een trede van de dalende: $1 + 4$, $2 + 3$,
$3 + 2$, $4 + 1$ --- altijd $5$). Twee keer de \omterm{def:g1:addition:def}{som} is dus $4 \times 5 = 20$
vakjes, en $1 + 2 + 3 + 4 = 10$: het plaatje \emph{is} vraag 7.

\textbf{14.} $1 + 2 + \dots + 12 = 78$ slagen in twaalf uur (vraag 12), en dus
$78 \times 2 = 156$ slagen op een hele dag.

\textbf{15.} Deel door $60$: $5\,050 = 60 \times 84 + 10$, dus $5\,050$
seconden is $84$ minuten en $10$ seconden. Deel de minuten door $60$:
$84 = 60 \times 1 + 24$, dus $84$ minuten is $1$ uur en $24$ minuten. Alles
samen: $1$ u $24$ min $10$ s --- iets \emph{minder} dan anderhalf uur ($1$ u
$30$ min).
\end{solution}
